\documentclass[12pt,letterpaper]{article}
\usepackage{graphicx,textcomp}
\usepackage{natbib}
\usepackage{setspace}
\usepackage{fullpage}
\usepackage{color}
\usepackage[reqno]{amsmath}
\usepackage{amsthm}
\usepackage{fancyvrb}
\usepackage{amssymb,enumerate}
\usepackage[all]{xy}
\usepackage{endnotes}
\usepackage{lscape}
\newtheorem{com}{Comment}
\usepackage{float}
\usepackage{hyperref}
\newtheorem{lem} {Lemma}
\newtheorem{prop}{Proposition}
\newtheorem{thm}{Theorem}
\newtheorem{defn}{Definition}
\newtheorem{cor}{Corollary}
\newtheorem{obs}{Observation}
\usepackage[compact]{titlesec}
\usepackage{dcolumn}
\usepackage{tikz}
\usetikzlibrary{arrows}
\usepackage{multirow}
\usepackage{xcolor}
\newcolumntype{.}{D{.}{.}{-1}}
\newcolumntype{d}[1]{D{.}{.}{#1}}
\definecolor{light-gray}{gray}{0.65}
\usepackage{url}
\usepackage{listings}
\usepackage{color}

\definecolor{codegreen}{rgb}{0,0.6,0}
\definecolor{codegray}{rgb}{0.5,0.5,0.5}
\definecolor{codepurple}{rgb}{0.58,0,0.82}
\definecolor{backcolour}{rgb}{0.95,0.95,0.92}

\lstdefinestyle{mystyle}{
	backgroundcolor=\color{backcolour},   
	commentstyle=\color{codegreen},
	keywordstyle=\color{magenta},
	numberstyle=\tiny\color{codegray},
	stringstyle=\color{codepurple},
	basicstyle=\footnotesize,
	breakatwhitespace=false,         
	breaklines=true,                 
	captionpos=b,                    
	keepspaces=true,                 
	numbers=left,                    
	numbersep=5pt,                  
	showspaces=false,                
	showstringspaces=false,
	showtabs=false,                  
	tabsize=2
}
\lstset{style=mystyle}
\newcommand{\Sref}[1]{Section~\ref{#1}}
\newtheorem{hyp}{Hypothesis}

\title{Problem Set 4}
\date{\today}
\author{Hanyu Li(25346841)}

\begin{document}
	\setlength{\parindent}{0pt}
	\setlength{\parskip}{0.8em}
	
	\maketitle
	
	\section*{Instructions}
	\begin{itemize}
		\item Please show your work! You may lose points by simply writing in the answer. If the problem requires you to execute commands in \texttt{R}, please include the code you used to get your answers. Please also include the \texttt{.R} file that contains your code. If you are not sure if work needs to be shown for a particular problem, please ask.
		\item Your homework should be submitted electronically on GitHub in \texttt{.pdf} form.
		\item This problem set is due before 23:59 on Friday April 10, 2026. No late assignments will be accepted.
	\end{itemize}
	
	\vspace{1.5cm}
	
	\section*{Question 1}
	
	We're interested in modeling the historical causes of child mortality. We have data from 26855 children born in Skellefteå, Sweden from 1850 to 1884. Using the "child" dataset in the \texttt{eha} library, fit a Cox Proportional Hazard model using mother's age and infant's gender as covariates. Present and interpret the output.
	
	\vspace{0.5cm}
	
	To estimate the historical causes of child mortality, we fit a Cox Proportional Hazard model. We used \texttt{Surv()} to create the survival object to predict the instantaneous risk (hazard) of child mortality given the mother's age and the infant's gender.
	
	\lstinputlisting[language=R, firstline=47, lastline=51]{PS04_HL.R}
	
	The regression results are as below:
	
	\begin{table}[H] \centering 
		\caption{Cox Proportional Hazard Model of Child Mortality (Hazard Ratios)} 
		\label{} 
		\begin{tabular}{@{\extracolsep{5pt}}lc} 
			\\[-1.8ex]\hline 
			\hline \\[-1.8ex] 
			& \multicolumn{1}{c}{\textit{Dependent variable:}} \\ 
			\cline{2-2} 
			\\[-1.8ex] & enter \\ 
			\hline \\[-1.8ex] 
			m.age & 1.008$^{***}$ \\ 
			& (0.002) \\ 
			& \\ 
			sexfemale & 0.921$^{***}$ \\ 
			& (0.027) \\ 
			& \\ 
			\hline \\[-1.8ex] 
			Observations & 26,574 \\ 
			R$^{2}$ & 0.001 \\ 
			Max. Possible R$^{2}$ & 0.986 \\ 
			Log Likelihood & $-$56,503.480 \\ 
			Wald Test & 22.520$^{***}$ (df = 2) \\ 
			LR Test & 22.518$^{***}$ (df = 2) \\ 
			Score (Logrank) Test & 22.530$^{***}$ (df = 2) \\ 
			\hline 
			\hline \\[-1.8ex] 
			\textit{Note:}  & \multicolumn{1}{r}{$^{*}$p$<$0.1; $^{**}$p$<$0.05; $^{***}$p$<$0.01} \\ 
		\end{tabular} 
	\end{table}  
	
	Based on output above, we can interpret the estimated hazard ratios ($\exp(\hat{\beta})$) as follows:
	
	\begin{itemize}
		\item Holding the infant's gender constant, a one-unit (year) increase in the mother's age leads to a change in the expected hazard of child mortality by a multiplicative factor of $\approx 1.008$ on average.
		\item Holding the mother's age constant, being a female infant versus a male infant (the baseline category) changes the expected hazard of child mortality by a multiplicative factor of $\approx 0.921$ on average.
	\end{itemize}
	
	Furthermore, we conducted a Likelihood Ratio Test (LRT) to assess our model quality. 
	
	The Full Model is specified as:
	
	$$h(t|X) = h_0(t) \exp(\beta_{\text{age}} \cdot \text{m.age} + \beta_{\text{sex}} \cdot \text{sex})$$
	
	We compared the full model against a reduced model where one specific covariate $k$ is omitted. For example, testing the effect of mother's age compares the full model to the reduced model: $h(t|X) = h_0(t) \exp(\beta_{\text{sex}} \cdot \text{sex})$. 
	
	The hypotheses for each covariate $k$ (where $k \in \{\text{m.age}, \text{sex}\}$) are defined as:
	
	\begin{itemize}
		\item $H_0: \beta_k = 0$
		\item $H_a: \beta_k \neq 0$
	\end{itemize}
	
	\lstinputlisting[language=R, firstline=62, lastline=63]{PS04_HL.R}
	
	\begin{verbatim}
		Single term deletions
		
		Model:
		Surv(enter, exit, event) ~ m.age + sex
		Df    AIC     LRT  Pr(>Chi)    
		<none>    113011                      
		m.age   1 113022 12.7946 0.0003476 ***
		sex     1 113018  9.4646 0.0020947 ** ---
		Signif. codes:  0 ‘***’ 0.001 ‘**’ 0.01 ‘*’ 0.05 ‘.’ 0.1 ‘ ’ 1 
	\end{verbatim}
	
	Based on the LRT results, p-values are nearly zero ($p \approx 0.0003$ for mother's age and $p \approx 0.002$ for infant's gender). Thus, we have enough evidence to reject the null hypothesis ($H_0: \beta_k = 0$) for both variables. This suggests that there is sufficient evidence to conclude that both the mother's age and the infant's gender reliably improve the explained variation of our survival model.
	
	\vspace{1.5cm}
	
	\section*{Question 2}
	
	We want to estimate how the amount of damage caused by a disaster relates to (1) whether disaster relief is provided and (2) the amount of relief that is provided conditional on a donation occurring when a disaster hits a given country. Estimate a Heckman selection model using the \texttt{disasters\_response} data on \href{https://raw.githubusercontent.com/ASDS-TCD/StatsII_2026/refs/heads/main/datasets/disaster_response.csv}{GitHub} in which \texttt{binContribution} is the outcome for the selection equation, and \texttt{originalContributionMillionUSDLogged} is the outcome for the second equation. Use \texttt{occurrences + deathsEM + normalizedDamageEMLogged} as your input variables for both equations. Present and interpret the output.
	
	\vspace{0.5cm}
	
	To estimate how the amount of damage caused by a disaster relates to (1) whether disaster relief is provided and (2) the amount of relief provided conditional on a donation occurring, we should account for potential selection bias. If we were to only run a standard regression on the subset of disasters that actually received donations, we would ignore the disasters that were not severe enough to receive international aid, which could lead to biased estimates. 
	
	To evaluate and correct for this bias, we construct a two-step Heckman selection model. 
	
	First, the Selection Equation (Probit) predicts the probability of a disaster receiving a donation:
	$$P(\text{binContribution} = 1) = \Phi(\gamma_0 + \gamma_1 \text{occurrences} + \gamma_2 \text{deathsEM} + \gamma_3 \text{normalizedDamageEMLogged})$$
	
	Second, the Outcome Equation (OLS) predicts the logged donation amount, including the Inverse Mills Ratio ($\lambda$) extracted from the first step to correct for selection bias:
	$$\text{originalContributionMillionUSDLogged} = \beta_0 + \beta_1 \text{occurrences} + \beta_2 \text{deathsEM} + \beta_3 \text{normalizedDamageEMLogged} + \beta_{\lambda}\lambda + \epsilon$$
	
	\lstinputlisting[language=R, firstline=75, lastline=81]{PS04_HL.R}
	
	The regression results for the outcome equation are presented in the table below:
	
	\begin{table}[H] \centering 
		\caption{Heckman Selection Model of Disaster Relief Actions (Outcome Equation)} 
		\label{} 
		\begin{tabular}{@{\extracolsep{5pt}}lc} 
			\\[-1.8ex]\hline 
			\hline \\[-1.8ex] 
			& \multicolumn{1}{c}{\textit{Dependent variable:}} \\ 
			\cline{2-2} 
			\\[-1.8ex] & originalContributionMillionUSDLogged \\ 
			\hline \\[-1.8ex] 
			occurrences & $-$0.086 \\ 
			& (0.055) \\ 
			& \\ 
			deathsEM & $-$0.033 \\ 
			& (0.026) \\ 
			& \\ 
			normalizedDamageEMLogged & $-$0.114$^{*}$ \\ 
			& (0.062) \\ 
			& \\ 
			Constant & 9.766 \\ 
			& (6.044) \\ 
			& \\ 
			\hline \\[-1.8ex] 
			Observations & 49,096 \\ 
			R$^{2}$ & 0.055 \\ 
			Adjusted R$^{2}$ & 0.053 \\ 
			$\rho$ & $-$1.053 \\ 
			Inverse Mills Ratio & $-$6.313$^{*}$  (3.420) \\ 
			\hline 
			\hline \\[-1.8ex] 
			\textit{Note:}  & \multicolumn{1}{r}{$^{*}$p$<$0.1; $^{**}$p$<$0.05; $^{***}$p$<$0.01} \\ 
		\end{tabular} 
	\end{table} 
	
	Based on the output, we observe that the \textbf{Inverse Mills Ratio} is only marginally significant at the 10\% level ($p < 0.1$) but is not statistically significant at the conventional 5\% alpha level ($p > 0.05$). This indicates that we do not have strong, reliable evidence of severe selection bias in this specific sample. Consequently, we can interpret the coefficients in the outcome equation straightforwardly, just as we would in a standard OLS regression for the observed data.
	
	The interpretations of the estimated OLS coefficients are as follows:
	\begin{itemize}
		\item \textbf{occurrences}: Holding all other variables constant, on average, a one-unit increase in the number of occurrences is associated with an estimated decrease of 0.086 in the logged original contribution (in million USD).
		\item \textbf{deathsEM}: Holding all other variables constant, on average, a one-unit increase in deaths is associated with an estimated decrease of 0.033 in the logged original contribution.
		\item \textbf{normalizedDamageEMLogged}: Holding all other variables constant, on average, a one-unit increase in the logged normalized damage is associated with an estimated decrease of 0.114 in the logged original contribution.
		\item \textbf{Constant (Intercept)}: When there are 0 disaster occurrences (\texttt{occurrences} = 0), 0 deaths (\texttt{deathsEM} = 0), and the raw normalized economic damage is 1 unit (\texttt{normalizedDamageEMLogged} = 0), the expected logged original contribution is, on average, 9.766.
	\end{itemize}
	
\end{document}